\section{Introduction}
A product of independent random switches can make each card uniform while leaving substantial dependence between cards. The Thorp shuffle illustrates this distinction particularly well. On $N=2^d$ positions, with $d\ge1$, one coordinate sweep visits the binary-coordinate directions in the fixed order $1,\ldots,d$. In each direction it independently swaps the two cards on every edge with probability $1/2$, leaving them in place otherwise. Every individual card is uniform after this sweep. The question is how many sweeps randomize the entire permutation.

We study that question through a four-row decomposition. Removing two coordinate directions leaves four smaller independent sweeps, one in each row of a $4$-by-$N/4$ array. The last two directions act within its columns. Each smaller sweep contracts according to the representations of its own row group, but restricting a representation of $S_N$ to the four row groups creates multiplicities. Those multiplicities accumulate as the decomposition is repeated. The main point is an inequality that makes their total cost strictly smaller than the available representation dimension.

For a function $f:\{1,2,3,4\}\to[0,\infty)$, write
\[
 \|f\|_p=\left(\frac14\sum_{j=1}^4 f(j)^p\right)^{1/p}.
\]
If $u$ denotes uniform measure on $S_4$, then
$\E_u\prod_i f_i(\pi(i))$ is the permanent of the matrix $(f_i(j))$ divided by $4!$. Carlen, Lieb and Loss proved that it is at most $\prod_i\|f_i\|_2$ and characterized equality \cite[Theorem~1.1]{cll2006}. For nonzero nonnegative functions on four points, equality requires that every function be constant. We will use both the bound and this equality statement.

\begin{samepage}
\begin{theorem}[A robust four-row inequality]
There are absolute constants $p_0\in(4/3,2)$ and $\varepsilon>0$ with the following property.\label{thm:permanent} Let $\nu$ be a probability measure on $S_4$ such that
\[
 \|\nu-u\|_{\TV}<\varepsilon,
 \qquad
 \nu\{\pi:\pi(i)=j\}=\frac14\quad(1\le i,j\le4).
\]
Then every four nonnegative functions on $\{1,2,3,4\}$ satisfy
\begin{equation}\label{eq:permanent}
 \E_\nu\prod_{i=1}^4f_i(\pi(i))\le\prod_{i=1}^4\|f_i\|_{p_0}.
\end{equation}
\end{theorem}
\end{samepage}

Here total variation is one half the sum of the absolute differences of the probability masses. The marginal assumption in the theorem is exact. Indeed, setting three functions equal to one in \eqref{eq:permanent}, and varying the fourth around the constant one, forces its coordinate marginal to be uniform. Thus this hypothesis describes a structural feature of the inequality, as well as the cancellation used in its proof.

Permanents with row or column $\ell^p$ constraints have a broader history. Carlen--Lieb--Loss discuss the below-two problem and Caputo's conjecture, and Samorodnitsky obtained further bounds in that range \cite{cll2006,Samorodnitsky2008}. Bristiel and Caputo subsequently proved the sharp permanent bound for every $p\ge1$ \cite[Corollary~1.14]{BristielCaputo2024}. In the present normalization, their result gives
\[
 \E_u\prod_i f_i(\pi(i))\le\prod_i\|f_i\|_p
 \quad\text{for every }p\ge p_c:=\frac{4\log4}{\log24}.
\]
The identity matrix shows that this uniform-law threshold is necessary.

The assertion needed here is stability under small perturbations of the permutation law that preserve every coordinate marginal. We prove this nonoptimal robust form directly from the exponent-two theorem: a strict quadratic gap treats functions near constants, while the equality classification and compactness treat all remaining functions. The local calculation explains the endpoint $4/3$ in Theorem~\ref{thm:permanent}; the final exponent is chosen closer to two and is not optimized. Thus a below-two bound for the uniform law is not itself the new ingredient in the recursion.

To state the application, let $\rho_\lambda$ be the irreducible unitary representation of $S_N$ indexed by a partition $\lambda\vdash N$, and write $D_\lambda$ for its dimension and $\lambda'$ for the conjugate partition. If $\mu_N$ is the law of one sweep, put
\[
 T_N(\lambda)=\sum_{g\in S_N}\mu_N(g)\rho_\lambda(g),
 \qquad B_N(\lambda)=T_N(\lambda)T_N(\lambda)^*.
\]
Traces are ordinary, unnormalized matrix traces. Unmarked logarithms are natural.

\begin{theorem}[A fixed moment with a degree saving]\label{thm:moment}
There exist an even integer $r\ge2$ and $\eta>0$, independent of $N$, such that for every dyadic $N$ and every $\lambda\vdash N$,
\begin{equation}\label{eq:moment}
 D_\lambda\Tr B_N(\lambda)^r\le D_\lambda^{1-\eta}.
\end{equation}
Consequently there is an absolute integer $q$ for which the total-variation distance of $q$ independent forward sweeps from uniform tends to zero as $N\to\infty$, uniformly over the starting deck.
\end{theorem}

A physical Thorp shuffle consists of independent fair switches on the pairs differing in one binary coordinate, followed by cyclic rotation of the coordinates. Tracking those deterministic rotations makes $d$ physical shuffles exactly one sweep. Theorem~\ref{thm:moment} therefore gives an $O(d)$ upper bound. The opposite order follows from the elementary support estimate
\begin{equation}\label{eq:support-intro}
 \|\mu_t-U_{S_N}\|_{\TV}\ge1-\frac{2^{tN/2}}{N!},
\end{equation}
where $\mu_t$ is the law after $t$ physical shuffles: there are only $2^{tN/2}$ coin strings. In particular, distance at most $1/4$ requires
$t\ge\lceil(2/N)\log_2(3N!/4)\rceil=2d-O(1)$.

The shuffle originates in Thorp's study of nonrandom shuffling and Faro \cite{Thorp1973}. For $N=2^d$, Morris obtained an $O(d^{44})$ upper bound \cite{Morris2008}, and Montenegro and Tetali improved it to $O(d^{29})$ \cite[Theorems~6.14--6.15]{MontenegroTetali2006}. Morris subsequently proved $O((\log N)^4)$ for every even deck size \cite{Morris2009} and $O(d^3)$ for dyadic decks \cite{Morris2013}. The present proof reaches logarithmic order through a four-row moment recursion.

The proof has three stages. Section~\ref{sec:permanent} proves the robust inequality and tensors it over columns. In Section~\ref{sec:recursion}, a positive trace comparison reduces the sweep moment to four smaller moments, multiplied by powers of squared restriction multiplicities. The power is
$\theta=1-1/p_0<1/2$: this strict improvement over one half is the feature the dimension comparison needs. Section~\ref{sec:multiplicities} bounds the multiplicities by removing the first few rows and columns of a Young diagram. Horizontal and vertical Pieri rules control the removed strips. For every fixed width $h$, the total number of boxes outside the first $h$ rows and columns of the four child diagrams is at most the corresponding number in the parent. We deduce this inequality from positive hook-Schur specialization, originating with Berele--Regev \cite{BereleRegev1987}; the formulas and tableau conventions used here are recorded in \cite{orellana-zabrocki2000,mason-niese2016}.

Finally, Section~\ref{sec:sparse} stops the recursion for diagrams close to a row or column. Section~\ref{sec:tree} sums the costs at every other scale, including the immediate children at which recursion has stopped, and compares the result with the hook-length formula. The strict inequality $\theta<1/2$ is exactly the margin that absorbs these costs. Section~\ref{sec:amplification} converts the resulting moment to mixing. This last step follows the finite-group Fourier approach of Diaconis--Shahshahani \cite{DiaconisShahshahani1981}, with matrix norms retained because a sweep is generally nonnormal.

The argument uses ordinary complex representation theory of symmetric groups: branching, induction and restriction, the Littlewood--Richardson and Pieri rules, and the hook-length formula \cite{James1978,Sagan2001,Stanley1999}. We give the additional multiplicity and dimension estimates in full. The positive trace comparison is the Araki--Lieb--Thirring inequality \cite{araki1990}. No shuffle moment or mixing theorem from another source is an input here.
